% --- Archon physics formalization source begin ---
% archon:physics
% archon:covers PhyXMiniProblems/problem_phyx_mini_0080.lean
% archon:source-report reports/phyx_mini/problem_phyx_mini_0080.source.json

\chapter{Physics problem phyx\_mini\_0080}
\label{ch:PhyXMiniProblems_problem_phyx_mini_0080}

\paragraph{Problem source.}
\#\# Physical scenario

In figure, first-order reflection from the reflection planes shown occurs when an x-ray beam of wavelength \$0.260 \textbackslash{}, \textbackslash{}text\{nm\}\$ makes an angle \$	heta = 63.8\textasciicircum{}\textbackslash{}circ\$ with the top face of the crystal.

\#\# Question

What is the unit cell size \$a\_0\$?

\#\# Answer choices

- A: A: 0.405nm

- B: B: 0.520nm

- C: C: 0.570nm

- D: D: 0.640nm

\#\# Figure caption (auxiliary; use the image as primary evidence)

The image depicts an atomic lattice structure typically associated with X-ray diffraction studies. Key elements include:

- **Green Dots**: Represent atoms arranged in a lattice.
- **Black Lines**: Connect the atoms, illustrating the geometric arrangement within the lattice.
- **Red Arrow and Text**: A red arrow labeled "X rays" shows the direction of incoming X-ray beams at an angle denoted by \textbackslash{}(\textbackslash{}theta\textbackslash{}) with respect to the lattice plane.
- **Angle (\textbackslash{}(\textbackslash{}theta\textbackslash{}))**: Indicates the angle of incidence of the X-rays on the lattice.
- **Distance Annotations**: Two distances labeled \textbackslash{}(a\_0\textbackslash{}). One is vertical, representing the spacing between layers of atoms, and the other is horizontal, marking the distance between atoms in the same plane.

This diagram is often used to explain Bragg's Law in the context of X-ray crystallography.

\#\# Dataset metadata

Wave Optics; Physical Model Grounding Reasoning, Multi-Formula Reasoning

\paragraph{Recorded answer/context.}
C: C: 0.570nm

\paragraph{Figure/image path.}
/root/proposal\_for\_physic/science-mango/phyx\_mini\_run/phyx\_data/test\_image/80.png

\paragraph{Formalization target.}
create a compiling Lean file with sorry bodies at `PhyXMiniProblems/problem\_phyx\_mini\_0080.lean`. The Lean declarations must preserve the physical quantities, dimensions or dimensional roles, figure labels, governing-law hypotheses, and final relation expressed by this problem.
Use Mathlib/Physlib names found through LeanExplore where available. If a physics API is missing, introduce faithful local abstractions rather than scalar placeholder aliases.

\begin{theorem}[Physics formalization target]
\label{thm:physics:phyx_mini_0080:target}
\lean{PhyXMiniProblems.ProblemPhyXMini0080.problem_phyx_mini_0080}
\uses{def:physics:phyx-mini-0080:phyxminiproblems-problemphyxmini0080-nanometersvalue, def:physics:phyx-mini-0080:phyxminiproblems-problemphyxmini0080-xraycrystalsetup, def:physics:phyx-mini-0080:phyxminiproblems-problemphyxmini0080-hasstatedreadouts, def:physics:phyx-mini-0080:phyxminiproblems-problemphyxmini0080-matchesshownlatticefigure, def:physics:phyx-mini-0080:phyxminiproblems-problemphyxmini0080-satisfiesdiagonalplanespacinglaw, def:physics:phyx-mini-0080:phyxminiproblems-problemphyxmini0080-satisfiesbraggreflectionlaw, def:physics:phyx-mini-0080:phyxminiproblems-problemphyxmini0080-hasphysicaldiffractionparameters, def:physics:phyx-mini-0080:phyxminiproblems-problemphyxmini0080-isclosestanswer}
The assigned autoformalize agent should translate this physics problem into Lean declarations in the covered file, with theorem and lemma proof bodies written as `by sorry`.
\end{theorem}
\begin{proof}
This is an autoformalization task, not a proof task. Produce statements that can later be proved by the physics prover without weakening the physical model.
\end{proof}

% --- Archon declaration topology begin ---
\section{Declaration topology}
\begin{definition}[Length Quantity]
  \label{def:physics:phyx-mini-0080:phyxminiproblems-problemphyxmini0080-lengthquantity}
  \lean{PhyXMiniProblems.ProblemPhyXMini0080.LengthQuantity}
  A signed physical length whose representations cohere across unit choices
\end{definition}
\begin{proof}
  This is the declared model definition of Length Quantity.
\end{proof}
\begin{definition}[length Value In]
  \label{def:physics:phyx-mini-0080:phyxminiproblems-problemphyxmini0080-lengthvaluein}
  \lean{PhyXMiniProblems.ProblemPhyXMini0080.lengthValueIn}
  \uses{def:physics:phyx-mini-0080:phyxminiproblems-problemphyxmini0080-lengthquantity}
  Read a physical length as a real scalar in the selected length unit
\end{definition}
\begin{proof}
  This is the declared model definition of length Value In.
\end{proof}
\begin{definition}[nanometers Value]
  \label{def:physics:phyx-mini-0080:phyxminiproblems-problemphyxmini0080-nanometersvalue}
  \lean{PhyXMiniProblems.ProblemPhyXMini0080.nanometersValue}
  \uses{def:physics:phyx-mini-0080:phyxminiproblems-problemphyxmini0080-lengthquantity, def:physics:phyx-mini-0080:phyxminiproblems-problemphyxmini0080-lengthvaluein}
  The numerical readout of a physical length in nanometres
\end{definition}
\begin{proof}
  This is the declared model definition of nanometers Value.
\end{proof}
\begin{definition}[degrees As Angle]
  \label{def:physics:phyx-mini-0080:phyxminiproblems-problemphyxmini0080-degreesasangle}
  \lean{PhyXMiniProblems.ProblemPhyXMini0080.degreesAsAngle}
  Interpret a scalar angle given in degrees as a Mathlib real angle
\end{definition}
\begin{proof}
  This is the declared model definition of degrees As Angle.
\end{proof}
\begin{definition}[Unit Cell Edge Label]
  \label{def:physics:phyx-mini-0080:phyxminiproblems-problemphyxmini0080-unitcelledgelabel}
  \lean{PhyXMiniProblems.ProblemPhyXMini0080.UnitCellEdgeLabel}
  The two a₀ arrows drawn in the square lattice cross-section
\end{definition}
\begin{proof}
  This is the declared model definition of Unit Cell Edge Label.
\end{proof}
\begin{definition}[Reflection Plane Family]
  \label{def:physics:phyx-mini-0080:phyxminiproblems-problemphyxmini0080-reflectionplanefamily}
  \lean{PhyXMiniProblems.ProblemPhyXMini0080.ReflectionPlaneFamily}
  The family of parallel diagonal reflection planes explicitly shown
\end{definition}
\begin{proof}
  This is the declared model definition of Reflection Plane Family.
\end{proof}
\begin{definition}[XRay Beam]
  \label{def:physics:phyx-mini-0080:phyxminiproblems-problemphyxmini0080-xraybeam}
  \lean{PhyXMiniProblems.ProblemPhyXMini0080.XRayBeam}
  \uses{def:physics:phyx-mini-0080:phyxminiproblems-problemphyxmini0080-lengthquantity}
  The monochromatic incident x-ray beam
\end{definition}
\begin{proof}
  This is the declared model definition of XRay Beam.
\end{proof}
\begin{definition}[Square Lattice Crystal]
  \label{def:physics:phyx-mini-0080:phyxminiproblems-problemphyxmini0080-squarelatticecrystal}
  \lean{PhyXMiniProblems.ProblemPhyXMini0080.SquareLatticeCrystal}
  \uses{def:physics:phyx-mini-0080:phyxminiproblems-problemphyxmini0080-lengthquantity}
  Crystal lengths relevant to the depicted square lattice cross-section
\end{definition}
\begin{proof}
  This is the declared model definition of Square Lattice Crystal.
\end{proof}
\begin{definition}[XRay Diffraction Figure]
  \label{def:physics:phyx-mini-0080:phyxminiproblems-problemphyxmini0080-xraydiffractionfigure}
  \lean{PhyXMiniProblems.ProblemPhyXMini0080.XRayDiffractionFigure}
  \uses{def:physics:phyx-mini-0080:phyxminiproblems-problemphyxmini0080-lengthquantity, def:physics:phyx-mini-0080:phyxminiproblems-problemphyxmini0080-unitcelledgelabel, def:physics:phyx-mini-0080:phyxminiproblems-problemphyxmini0080-reflectionplanefamily}
  Geometric quantities and labels read from the primary figure
\end{definition}
\begin{proof}
  This is the declared model definition of XRay Diffraction Figure.
\end{proof}
\begin{definition}[Reflection Observation]
  \label{def:physics:phyx-mini-0080:phyxminiproblems-problemphyxmini0080-reflectionobservation}
  \lean{PhyXMiniProblems.ProblemPhyXMini0080.ReflectionObservation}
  The observed reflected order and the two angles used in Bragg's law
\end{definition}
\begin{proof}
  This is the declared model definition of Reflection Observation.
\end{proof}
\begin{definition}[XRay Crystal Setup]
  \label{def:physics:phyx-mini-0080:phyxminiproblems-problemphyxmini0080-xraycrystalsetup}
  \lean{PhyXMiniProblems.ProblemPhyXMini0080.XRayCrystalSetup}
  \uses{def:physics:phyx-mini-0080:phyxminiproblems-problemphyxmini0080-xraybeam, def:physics:phyx-mini-0080:phyxminiproblems-problemphyxmini0080-squarelatticecrystal, def:physics:phyx-mini-0080:phyxminiproblems-problemphyxmini0080-xraydiffractionfigure, def:physics:phyx-mini-0080:phyxminiproblems-problemphyxmini0080-reflectionobservation}
  Complete physical setup for the x-ray reflection observation
\end{definition}
\begin{proof}
  This is the declared model definition of XRay Crystal Setup.
\end{proof}
\begin{definition}[Has Stated Readouts]
  \label{def:physics:phyx-mini-0080:phyxminiproblems-problemphyxmini0080-hasstatedreadouts}
  \lean{PhyXMiniProblems.ProblemPhyXMini0080.HasStatedReadouts}
  \uses{def:physics:phyx-mini-0080:phyxminiproblems-problemphyxmini0080-nanometersvalue, def:physics:phyx-mini-0080:phyxminiproblems-problemphyxmini0080-degreesasangle, def:physics:phyx-mini-0080:phyxminiproblems-problemphyxmini0080-xraycrystalsetup}
  This structure records the Has Stated Readouts component of the problem-specific physical model
\end{definition}
\begin{proof}
  This is the declared model definition of Has Stated Readouts.
\end{proof}
\begin{definition}[Matches Shown Lattice Figure]
  \label{def:physics:phyx-mini-0080:phyxminiproblems-problemphyxmini0080-matchesshownlatticefigure}
  \lean{PhyXMiniProblems.ProblemPhyXMini0080.MatchesShownLatticeFigure}
  \uses{def:physics:phyx-mini-0080:phyxminiproblems-problemphyxmini0080-degreesasangle, def:physics:phyx-mini-0080:phyxminiproblems-problemphyxmini0080-xraycrystalsetup}
  This structure records the Matches Shown Lattice Figure component of the problem-specific physical model
\end{definition}
\begin{proof}
  This is the declared model definition of Matches Shown Lattice Figure.
\end{proof}
\begin{definition}[Satisfies Diagonal Plane Spacing Law]
  \label{def:physics:phyx-mini-0080:phyxminiproblems-problemphyxmini0080-satisfiesdiagonalplanespacinglaw}
  \lean{PhyXMiniProblems.ProblemPhyXMini0080.SatisfiesDiagonalPlaneSpacingLaw}
  \uses{def:physics:phyx-mini-0080:phyxminiproblems-problemphyxmini0080-xraycrystalsetup}
  This def records the Satisfies Diagonal Plane Spacing Law component of the problem-specific physical model
\end{definition}
\begin{proof}
  This is the declared model definition of Satisfies Diagonal Plane Spacing Law.
\end{proof}
\begin{definition}[Satisfies Bragg Reflection Law]
  \label{def:physics:phyx-mini-0080:phyxminiproblems-problemphyxmini0080-satisfiesbraggreflectionlaw}
  \lean{PhyXMiniProblems.ProblemPhyXMini0080.SatisfiesBraggReflectionLaw}
  \uses{def:physics:phyx-mini-0080:phyxminiproblems-problemphyxmini0080-xraycrystalsetup}
  This def records the Satisfies Bragg Reflection Law component of the problem-specific physical model
\end{definition}
\begin{proof}
  This is the declared model definition of Satisfies Bragg Reflection Law.
\end{proof}
\begin{definition}[Has Physical Diffraction Parameters]
  \label{def:physics:phyx-mini-0080:phyxminiproblems-problemphyxmini0080-hasphysicaldiffractionparameters}
  \lean{PhyXMiniProblems.ProblemPhyXMini0080.HasPhysicalDiffractionParameters}
  \uses{def:physics:phyx-mini-0080:phyxminiproblems-problemphyxmini0080-nanometersvalue, def:physics:phyx-mini-0080:phyxminiproblems-problemphyxmini0080-xraycrystalsetup}
  Positivity and physical-branch conditions for the depicted observation
\end{definition}
\begin{proof}
  This is the declared model definition of Has Physical Diffraction Parameters.
\end{proof}
\begin{definition}[Answer Choice]
  \label{def:physics:phyx-mini-0080:phyxminiproblems-problemphyxmini0080-answerchoice}
  \lean{PhyXMiniProblems.ProblemPhyXMini0080.AnswerChoice}
  Labels of the four displayed multiple-choice answers
\end{definition}
\begin{proof}
  This is the declared model definition of Answer Choice.
\end{proof}
\begin{definition}[answer Unit Cell Size Nanometers]
  \label{def:physics:phyx-mini-0080:phyxminiproblems-problemphyxmini0080-answerunitcellsizenanometers}
  \lean{PhyXMiniProblems.ProblemPhyXMini0080.answerUnitCellSizeNanometers}
  \uses{def:physics:phyx-mini-0080:phyxminiproblems-problemphyxmini0080-answerchoice}
  Unit-cell sizes printed beside the answer choices, in nanometres
\end{definition}
\begin{proof}
  This is the declared model definition of answer Unit Cell Size Nanometers.
\end{proof}
\begin{definition}[Is Closest Answer]
  \label{def:physics:phyx-mini-0080:phyxminiproblems-problemphyxmini0080-isclosestanswer}
  \lean{PhyXMiniProblems.ProblemPhyXMini0080.IsClosestAnswer}
  \uses{def:physics:phyx-mini-0080:phyxminiproblems-problemphyxmini0080-lengthquantity, def:physics:phyx-mini-0080:phyxminiproblems-problemphyxmini0080-nanometersvalue, def:physics:phyx-mini-0080:phyxminiproblems-problemphyxmini0080-answerchoice, def:physics:phyx-mini-0080:phyxminiproblems-problemphyxmini0080-answerunitcellsizenanometers}
  A choice is strictly nearer to a length readout than every other choice
\end{definition}
\begin{proof}
  This is the declared model definition of Is Closest Answer.
\end{proof}
\begin{lemma}[bragg glancing angle eq eighteen point eight degrees]
  \label{lem:physics:phyx-mini-0080:phyxminiproblems-problemphyxmini0080-bragg-glancing-angle-eq-eighteen-point-eight-degrees}
  \lean{PhyXMiniProblems.ProblemPhyXMini0080.bragg_glancing_angle_eq_eighteen_point_eight_degrees}
  \uses{def:physics:phyx-mini-0080:phyxminiproblems-problemphyxmini0080-degreesasangle, def:physics:phyx-mini-0080:phyxminiproblems-problemphyxmini0080-xraycrystalsetup, def:physics:phyx-mini-0080:phyxminiproblems-problemphyxmini0080-hasstatedreadouts, def:physics:phyx-mini-0080:phyxminiproblems-problemphyxmini0080-matchesshownlatticefigure}
  This lemma records the bragg glancing angle eq eighteen point eight degrees component of the problem-specific physical model
\end{lemma}
\begin{proof}
  The conclusion follows from the governing relations and figure readouts named in the statement of bragg glancing angle eq eighteen point eight degrees.
\end{proof}
\begin{lemma}[unit cell size nanometers formula]
  \label{lem:physics:phyx-mini-0080:phyxminiproblems-problemphyxmini0080-unit-cell-size-nanometers-formula}
  \lean{PhyXMiniProblems.ProblemPhyXMini0080.unit_cell_size_nanometers_formula}
  \uses{def:physics:phyx-mini-0080:phyxminiproblems-problemphyxmini0080-nanometersvalue, def:physics:phyx-mini-0080:phyxminiproblems-problemphyxmini0080-xraycrystalsetup, def:physics:phyx-mini-0080:phyxminiproblems-problemphyxmini0080-hasstatedreadouts, def:physics:phyx-mini-0080:phyxminiproblems-problemphyxmini0080-satisfiesdiagonalplanespacinglaw, def:physics:phyx-mini-0080:phyxminiproblems-problemphyxmini0080-satisfiesbraggreflectionlaw, def:physics:phyx-mini-0080:phyxminiproblems-problemphyxmini0080-hasphysicaldiffractionparameters}
  This lemma records the unit cell size nanometers formula component of the problem-specific physical model
\end{lemma}
\begin{proof}
  The conclusion follows from the governing relations and figure readouts named in the statement of unit cell size nanometers formula.
\end{proof}
% --- Archon declaration topology end ---
% --- Archon physics formalization source end ---
